\section{Proof of the Global SPD Action Bundle}
\label{app:action-bundle}

\begin{proof}[Proof of \cref{thm:action-bundle}]
We proceed in seven steps.

\paragraph{Step 1: coordinates on the action base.}
Put \(Y=BC^{-1}\).  Since \(Q=[U,V]\) is orthogonal,
\[
  Y=UM+VN,
  \qquad M=U^\top Y,
  \qquad N=V^\top Y.
\]
Moreover,
\[
  A^\top B=C^\top M C,
\]
so \(A^\top B\in\SPD^m\) exactly when \(M\in\SPD^m\).  The map
\(B\mapsto(M,N)\) has the analytic inverse
\begin{equation}
  B=(UM+VN)C.
  \label{eq:appendix-base-inverse}
\end{equation}
Thus \(\Bact\cong\SPD^m\times\mathbb R^{n\times m}\).

\paragraph{Step 2: every SPD completion.}
Suppose \(P\in\SPD^d\) satisfies \(PA=B\).  Because \(A=UC\), this is
equivalent to \(PU=Y\).  In the \(Q\) coordinates, symmetry fixes the first
block row and column:
\[
  \widehat P:=Q^\top P Q
  =\begin{pmatrix}M&N^\top\\N&D\end{pmatrix}.
\]
The Schur-complement criterion gives
\[
  \widehat P\succ0
  \quad\Longleftrightarrow\quad
  S:=D-NM^{-1}N^\top\succ0.
\]
With \(T=NM^{-1}\), every positive solution and only such a solution is
\begin{equation}
  \widehat P
  =L\begin{pmatrix}M&0\\0&S\end{pmatrix}L^\top.
  \label{eq:appendix-all-completions}
\end{equation}
Since \(\det L=1\),
\[
  \det P=\det M\det S.
\]
The determinant-one constraint is therefore equivalent to
\[
  S=\rho\Sigma,
  \qquad \rho=(\det M)^{-1/n},
  \qquad \det\Sigma=1,
\]
which is exactly \cref{eq:action-bundle-map}.  Conversely, direct block
multiplication and \cref{eq:appendix-base-inverse} give
\[
  \Phi_A(B,\Sigma)A
  =Q\begin{pmatrix}M\\N\end{pmatrix}C=B.
\]

\paragraph{Step 3: analytic inverse.}
Given \(P\in\Pdet\), recover \(B=PA\), then \(M(B),N(B)\), and the unique
Schur complement
\[
  S=D-NM^{-1}N^\top.
\]
The unique determinant-one hidden coordinate is
\begin{equation}
  \Sigma=(\det M)^{1/n}S.
  \label{eq:appendix-hidden-inverse}
\end{equation}
All operations are analytic on their SPD domains.  Thus \(\Phi_A\) is a
global analytic diffeomorphism and \(\pi_A\) is projection onto its first
factor.

\paragraph{Step 4: fiber geometry.}
Fix \(B\).  Before congruence by \(QL\), the fiber is
\[
  \mathcal C_B
  =\{\operatorname{diag}(M,\rho\Sigma):\Sigma\in\mathscr P_n\}.
\]
The ambient AIRM geodesic between two such points is
\[
  \operatorname{diag}(M,\rho\Sigma_0)
  \#_t
  \operatorname{diag}(M,\rho\Sigma_1)
  =\operatorname{diag}(M,\rho(\Sigma_0\#_t\Sigma_1)).
\]
Because log determinant is affine along AIRM geodesics,
\(\det(\Sigma_0\#_t\Sigma_1)=1\).  Hence the fiber is totally geodesic.
It is closed, and block distance splitting gives
\[
  \dai(\operatorname{diag}(M,\rho\Sigma_0),
       \operatorname{diag}(M,\rho\Sigma_1))
  =\dai(\Sigma_0,\Sigma_1).
\]
Congruence by \(QL\) is an AIRM isometry, proving all fiber claims.

\paragraph{Step 5: canonical projection and analyticity.}
The determinant-one SPD slice is a finite-dimensional Hadamard manifold.  A
nonempty closed geodesically convex set has a unique metric projection, so
\cref{eq:canonical-controller} defines a global section.

In the global coordinates define
\[
  F(B,\Sigma)=\frac12\fro{\log\Phi_A(B,\Sigma)}^2.
\]
This function is analytic.  For fixed \(B\), its hidden variable parameterizes
an isometric totally geodesic fiber.  Squared distance to \(I\), restricted to
that fiber, is strongly geodesically convex.  At the unique critical point the
vertical Hessian is positive definite.  The analytic implicit-function theorem
therefore gives a local analytic minimizer \(\Sigma_\star(B)\); global
uniqueness makes these local sections agree.  Hence
\[
  s_A(B)=\Phi_A(B,\Sigma_\star(B))
\]
and \(\mathfrak C_A(B)^2=2F(B,\Sigma_\star(B))\) are analytic.

\paragraph{Step 6: equivariance.}
For \(D\in\operatorname{GL}_m\),
\[
  P(AD)=BD\quad\Longleftrightarrow\quad PA=B,
\]
so the two action fibers are identical and have the same unique projection.
For \(R\in O(d)\),
\[
  PA=B
  \quad\Longleftrightarrow\quad
  (RPR^\top)(RA)=RB.
\]
Orthogonal congruence preserves determinant and distance to \(I\), so
uniqueness gives \cref{eq:controller-equivariance}.

\paragraph{Step 7: normal equation and certificate.}
Congruence invariance reduces the hidden objective to
\[
  f_B(\Sigma)=\frac12\dai(R_0,\mathcal Q_\Sigma)^2.
\]
We use the convention that a function is \(1\)-strongly geodesically convex
when every constant-speed geodesic \(\gamma_t\) satisfies
\begin{equation}
  f(\gamma_t)
  \le (1-t)f(\gamma_0)+t f(\gamma_1)
  -\frac12t(1-t)d(\gamma_0,\gamma_1)^2.
  \label{eq:cat0-strong-convexity-convention}
\end{equation}
On a CAT(0) space this inequality holds for
\(f(\cdot)=\tfrac12d(\cdot,R_0)^2\).  The AIRM SPD manifold is Hadamard, and
restriction to the isometric totally geodesic determinant-one fiber preserves
the same constant \cite{bhatia2007positive,lim2004best}.
A whitened tangent to the determinant-one hidden block has the form
\[
  Z=\begin{pmatrix}0&0\\0&Z_{22}\end{pmatrix},
  \qquad Z_{22}=Z_{22}^\top,
  \qquad \tr Z_{22}=0.
\]
The whitened negative gradient is the tangent projection of
\(\mathcal R(\Sigma)\) in \cref{eq:action-residual}.  Orthogonality to every
trace-free \(Z_{22}\) is exactly
\(\mathcal R_{22}(\Sigma)^\circ=0\).  Strong convexity makes this condition
necessary and sufficient.

Let \(D=\dai(\Sigma,\Sigma_\star)\), let
\(\xi=\operatorname{Log}_\Sigma(\Sigma_\star)\), and let \(V(\Sigma)\) be the restricted
negative gradient.  The first-order form of
\cref{eq:cat0-strong-convexity-convention} gives
\[
  f_B(\Sigma)-f_B(\Sigma_\star)
  \le \langle V(\Sigma),\xi\rangle-\frac12D^2
  \le\norm{V(\Sigma)}D-\frac12D^2,
\]
while the same inequality based at the minimizer, whose gradient vanishes,
gives
\[
  f_B(\Sigma)-f_B(\Sigma_\star)\ge\frac12D^2.
\]
Adding the two bounds yields
\[
  D^2\le\norm{V(\Sigma)}D,
\]
and hence, also for \(D=0\),
\[
  D\le\norm{V(\Sigma)}
  =\fro{\mathcal R_{22}(\Sigma)^\circ}.
\]
Finally,
\[
  f_B(\Sigma)-f_B(\Sigma_\star)
  \le \sup_{D\ge0}
  \left(\norm{V(\Sigma)}D-\frac12D^2\right)
  =\frac12\norm{V(\Sigma)}^2.
\]
This proves the exact constants in
\cref{eq:action-distance-certificate,eq:action-value-certificate}.
\end{proof}

\subsection{Proof of active reduction}

Let \(J\) be the orthogonal reflection that equals \(+I\) on \(W\) and
\(-I\) on \(W^\perp\).  Since the columns of \(A\) and \(B\) lie in \(W\),
\[
  JA=A,
  \qquad JB=B.
\]
If \(P\in\Pfiber(B)\), then \((JPJ)A=B\), and \(JPJ\) has the same
determinant and distance from \(I\) as \(P\).  Uniqueness of the canonical
controller yields
\[
  Js_A(B)J=s_A(B).
\]
Hence the cross blocks between \(W\) and \(W^\perp\) vanish:
\(s_A(B)=H_\star\oplus C_\star\).

For fixed feasible \(H\), the inactive log-eigenvalues \(z_i\) satisfy
\[
  \sum_{i=1}^qz_i=-\log\det H.
\]
Cauchy--Schwarz gives
\[
  \fro{\log C}^2=\sum_i z_i^2
  \ge\frac{(\log\det H)^2}{q},
\]
with equality exactly when all \(z_i\) are equal.  Thus
\[
  C_\star=(\det H_\star)^{-1/q}I_q,
\]
and substitution yields \cref{eq:active-controller-problem}.  Uniqueness
follows from uniqueness of \(s_A(B)\).
